% !TeX program = xelatex
\documentclass[11pt,a4paper]{article}
\usepackage[margin=23mm]{geometry}
\usepackage{fontspec}
\setmainfont{Latin Modern Roman}
\usepackage{amsmath,amssymb,mathtools}
\usepackage{xcolor,enumitem,booktabs,longtable,array,xurl,hyperref}
\definecolor{muted}{HTML}{53636C}
\hypersetup{colorlinks=true,urlcolor=blue,linkcolor=blue}
\setlength{\parindent}{0pt}
\setlength{\parskip}{5pt}
\newcommand{\R}{\mathbb R}
\newcommand{\E}{\mathbb E}
\newcommand{\N}{\mathbb N}
\newcommand{\ind}{\mathbf 1}
\newcommand{\Var}{\operatorname{Var}}
\newcommand{\Cov}{\operatorname{Cov}}
\newcommand{\supp}{\operatorname{supp}}
\DeclareMathOperator*{\argmin}{arg\,min}
\DeclareMathOperator*{\argmax}{arg\,max}
\newcommand{\entrymeta}[3]{{\sffamily\small\color{muted}\textbf{\texttt{#1}}\quad|\quad #2\par #3\par}}
\newcommand{\Problemref}[1]{\texttt{#1}}
\newcommand{\JELref}[2]{\texttt{#2} (JEL \texttt{#1})}
\begin{document}
\section*{Separate identification of beliefs and preferences}
\textbf{Problem ID:} \texttt{D81-16}\quad \textbf{JEL:} \texttt{D81}\par
% BEGIN ECONOMICS STATEMENT
\label{problem:OP-0587}
\label{beyond:DT-02}
\entrymeta{DT-02}{Decision theory and identification}{Research agenda formalization}
\subsection*{Definitions and assumptions}
This is a characterization agenda across declared model classes and designs, rather than an injectivity assertion for one unspecified experiment. For each class, let $m$ be a positive integer, $\Omega=\{1,\ldots,m\}$, and $\Delta_m=\{p\in\R^m:p_s\geq0,\ \sum_{s=1}^m p_s=1\}$. Thus $m=1$ and boundary priors are allowed unless a proposed class explicitly restricts them. Specify a parameter domain $\Theta$, an ordered consequence domain $C$ containing $c_0<c_1$, and a family of real-valued utilities $u_\theta:C\to\R$, $\theta\in\Theta$. State whether increasing means weakly or strictly increasing. Require a positive utility gap at the anchors before normalization, and normalize $u_\theta(c_0)=0$, $u_\theta(c_1)=1$ to remove positive-affine utility equivalence. This does not remove duplicate parameter labels: identification of $\theta$ itself must distinguish such labels, whereas identification of normalized utility treats labels with the same function on $C$ as equivalent. Declare which target is used. Specify measurable environment and action spaces, a supported environment set $Z$, nonempty menus $A(z)$, and a known consequence map $c(a,s,z)\in C$ for $z\in Z$, $a\in A(z)$ and $s\in\Omega$. For the maximizing-choice model below, require attainment and the measurability needed by the declared observation mechanism; finite menus are sufficient for attainment but are not required. Under expected utility, the choice correspondence is
\[
 \mathcal C_{p,\theta}(z)=\argmax_{a\in A(z)}
 \sum_{s=1}^{m}p_su_\theta(c(a,s,z)).
\]
The displayed formula uses one fixed pair $(p,\theta)$ across environments. Each design must specify the decision maker's information before choice and the restrictions under which menu or payoff variation leaves this pair invariant. If environments convey state information and beliefs update, declare the information and updating model and replace $p_s$ in the evaluation by the resulting conditional state probabilities. If observations mix heterogeneous types, declare their distribution and target; a mixture law is not indexed by a single pair alone.

Declare the observational unit, the observable variables and their measurable response space, and whether the full maximizing correspondence or a selected action is observed. Specify any tie-selection, choice-error and belief-reporting mechanisms, including the response format, events elicited, rounding and the exact meaning of any truthful-reporting assumption. These mechanisms must determine a well-defined observation law. Known mechanisms are fixed inputs; unknown mechanisms have nuisance parameters $\eta$ with a declared admissible domain and restrictions. Write $M$ for a structural model in the resulting class $\mathfrak M$, including $(p,\theta,\eta)$ where applicable, and $\mathcal O_M$ for its population observation law. Declare either a fixed design law $\mu$ on $Z$, with conditional-law equality understood $\mu$-almost everywhere, or an experiment observing the specified response law at every $z\in Z$, with equality pointwise on $Z$. Here population law denotes the theoretical distribution of the declared observations, not necessarily a heterogeneous population. When mechanisms are known and the fixed-pair formula applies, write $\mathcal O(p,\theta)$ for this law. Two structural models are observationally equivalent when their declared observation laws agree.
\subsection*{Formal question}
For explicitly declared classes $\mathfrak M$ and designs of the preceding kind, characterize primitive payoff and menu variation, truthful elicitation assumptions, utility restrictions and exclusion restrictions that make $\mathcal O$ injective in $(p,\theta)$ after normalization when the observation law is indexed by that pair alone. With nuisance parameters, separate identification of the declared belief--preference target $T(M)$ means
\[
 \mathcal O_{M_1}=\mathcal O_{M_2}\ \Longrightarrow\ T(M_1)=T(M_2)
 \qquad\text{for all }M_1,M_2\in\mathfrak M,
\]
where $T(M)$ is $(p,\theta)$ or $(p,u_\theta)$ as declared, with the corresponding target explicitly specified for updating or heterogeneous-type models. This quantification compares all admissible nuisance values and does not require the nuisance itself to be identified. At a compatible observation law $P$, characterize the sharp class $\{M\in\mathfrak M:\mathcal O_M=P\}$ and its target image. For every policy-relevant functional under study, specify a map $H:\mathfrak M\to\mathcal H$ and any policy or intervention assumptions it uses; its sharp identified set is $\{H(M):M\in\mathfrak M,\ \mathcal O_M=P\}$, and it is identified at $P$ exactly when this nonempty set is a singleton. Distinguish this law-specific question from global identification over the whole class. An empty compatible class denotes model incompatibility. The requested characterizations should give substantive conditions in primitives or informative descriptions of these classes and images beyond restating observational equivalence. The task concerns population identification; finite-sample estimation, computational recognition and minimal experiment size are separate questions unless explicitly added for a proposed class.

Extend the question to a declared ambiguity-sensitive criterion $V(a;\Pi,\theta)$. Each such extension must specify the admissible nonempty prior sets $\Pi\subseteq\Delta_m$, the criterion's evaluation formula and any additional parameters, utility restrictions and normalizations, and the observation mechanism. Declare whether the belief target is the raw prior set or an equivalence class under a specified representation equivalence; distinct representing sets need not yield distinct evaluations. Apply the same global and law-specific identification questions to that declared target, keeping the prior-set representation and preference parameters distinct. No particular ambiguity criterion or representation normalization is imposed by this agenda.
\subsection*{Scope and status}
An exclusion restriction must specify which experimental or environmental variation changes consequences while leaving beliefs or preferences invariant; merely naming an instrument is insufficient. Objective lotteries can inform utility only under a stated model of their evaluation. Finite deterministic choices generally give inequalities rather than exact recovery of unrestricted continuous parameters. The task is to identify informative designs or model classes, not to assert that unrestricted beliefs and preferences are always separately identifiable. The source motivates this confounding problem but does not pose the displayed map-injectivity question as a numbered conjecture.
\subsection*{References for the source of the problem}
Charles F.~Manski, \emph{Survey Measurement of Probabilistic Macroeconomic Expectations: Progress and Promise}, \emph{NBER Macroeconomics Annual} 32 (2018), 411--471, Section~II.A, ``Economic Research on Choice under Uncertainty, without Expectations Data,'' and Section~VII.C, ``Confounding Beliefs and Preferences.'' \url{https://www.journals.uchicago.edu/doi/full/10.1086/696061}. NBER Working Paper~23418 (2017), \url{https://www.nber.org/papers/w23418}.

\subsection*{Common conventions: Source mathematical conventions}

\textbf{Well-defined objects.}
All probability spaces and measurable variables are specified or inherited
from a local statistical model. Expectations used as finite targets require
the stated integrability. A decision rule or estimator is measurable with
respect to its available information; randomized rules may use an auxiliary
independent random variable. Norms without a subscript are Euclidean in
finite-dimensional spaces. An optimizer is requested only when attainment
is assured; otherwise the question concerns the optimal value and
$\varepsilon$-optimal admissible rules for every $\varepsilon>0$.

\textbf{Identification.}
A structural model $m\in\mathcal M$ implies an observable law
$P_m^{\rm obs}$ and target $\theta(m)$. The target is point identified on
$\mathcal M$ if
\[
 P_{m_1}^{\rm obs}=P_{m_2}^{\rm obs}
 \quad\Longrightarrow\quad\theta(m_1)=\theta(m_2)
 \qquad(m_1,m_2\in\mathcal M).
\]
When this fails, the sharp identified set is
\[
 \Theta(P;\mathcal M)=\{\theta(m):m\in\mathcal M,
                                  P_m^{\rm obs}=P\}.
\]
Specifying an intervention or estimand does not establish identification.
Positivity, exclusion, causal sufficiency, measurement restrictions, and
transport assumptions are substantive local inputs, never automatic
consequences of notation.

\textbf{Minimax risk and nontrivial inference.}
For a statistical experiment with data law $P^{(n)}$, target $\theta(P)$,
and nonnegative loss $\ell$, define
\[
 \mathcal R_n(\mathcal P,\ell)
   =\inf_{\widehat\theta_n}\sup_{P\in\mathcal P}
      \E_{P^{(n)}}\ell(\widehat\theta_n,\theta(P)).
\]
The infimum is over measurable estimators. A sharp rate is a sequence
$r_n$ for which $0<c\leq C<\infty$, independent of $n$, satisfies
$c r_n\leq\mathcal R_n\leq C r_n$ for all sufficiently large $n$.
In suitable regular identifiable parametric experiments, under the conditions
needed for matching risk bounds, squared-error risk is of order $n^{-1}$;
the corresponding root-mean-squared error is of order $n^{-1/2}$.
Parametric structure alone does not guarantee these rates, and no estimation
rate is asserted for the present identification question. These different conventions
are distinguished locally. A uniform asymptotic confidence region $C_n$
must satisfy
\[
 \liminf_{n\to\infty}\inf_{P\in\mathcal P}
     \Pr_{P^{(n)}}\{\theta(P)\in C_n\}\geq 1-\alpha,
 \qquad 0<\alpha<1,
\]
together with an informative diameter or expected-width requirement.
Coverage alone permits the vacuous whole parameter space. Testing questions
similarly impose both size control and power against a declared separated
alternative class.

\textbf{Binary causal baseline.}
When an entry explicitly invokes this baseline, one observes independent
copies of $O=(X,W,Y)$, with $W\in\{0,1\}$ and potential outcomes $Y(0),Y(1)$.
Assume consistency $Y=Y(W)$, no interference, conditional exchangeability
$(Y(0),Y(1))\perp W\mid X$, and overlap
$\epsilon\leq e(X)\leq1-\epsilon$ almost surely for a fixed
$0<\epsilon<1/2$, where
\[
 e(x)=\Pr(W=1\mid X=x),\qquad
 m_w(x)=\E[Y\mid W=w,X=x],\qquad
 \tau(x)=m_1(x)-m_0(x).
\]
These assumptions identify conditional mean effects almost everywhere
under the covariate law. Pointwise or uniform effects require appropriate
regularity and a specified support. Continuous treatment, longitudinal
data, adaptive experiments, and network interference require their own
assumptions in place of this baseline. Bounded outcomes, densities, and
smoothness are additional assumptions only where stated.

\textbf{Function classes.}
On $[0,1]^d$, $\mathcal H_d^s(L)$ is the isotropic Hölder ball of radius
$L>0$: put $k=\lceil s\rceil-1$ and $\gamma=s-k\in(0,1]$, bound derivatives
through order $k$, and require the order-$k$ derivatives to be
$\gamma$-Hölder with constant at most $L$. Thus integer orders use a
Lipschitz final derivative convention. The class
$\operatorname{BV}_d(L)$ consists of functions $f\in L^1((0,1)^d)$ whose
distributional derivative is a finite vector measure and for which
$\|f\|_{L^1}+|Df|((0,1)^d)\leq L$.
An $L^\infty$ bound or a particular pointwise version is imposed separately
when needed. Sparse and neural-network classes require a declared
dictionary or architecture and coefficient bounds.

An anisotropic H\"older class with declared block smoothness indices means
coordinatewise H\"older regularity in each block, uniformly in the other
blocks, with all corresponding derivative and H\"older bounds at most
the stated radius. Derivative orders and the integer-order convention in
each block are those just given. Mixed-derivative bounds are additional
restrictions only when explicitly stated.

\textbf{Decision and computational questions.}
Policy regret compares admissible feasible policies under the same law and
constraints. In computational entries, finite data and thresholds have
rational encodings; running time is measured in their total bit length.
Real-valued equilibrium search uses the explicitly declared algebraic or
FIXP encoding. An approximation ratio for maximization is written
$\operatorname{OPT}/\operatorname{ALG}$ and for minimization as
$\operatorname{ALG}/\operatorname{OPT}$, with zero optima and infeasible
instances handled locally. Historical approximation bounds are not
automatically current bounds.
% END ECONOMICS STATEMENT
\end{document}
